\documentclass[11pt]{article}

\usepackage[margin=1in]{geometry}
\usepackage[T2A]{fontenc}
\usepackage[utf8]{inputenc}
\usepackage[english,russian]{babel}
\usepackage{amsmath}
\usepackage{amssymb}
\usepackage{booktabs}
\usepackage{hyperref}

\title{Конечнополевая edge-инференция с ядрами GF(137) и компактная эмуляция виртуальных вентилей}
\author{И. Михайлов\\
\small Независимый исследователь\\
\small \href{https://doi.org/10.5281/zenodo.20436771}{doi:10.5281/zenodo.20436771}}
\date{Май 2026}

\begin{document}
\maketitle

\begin{abstract}
Edge-инференция часто ограничена не только числом параметров модели, но и
пропускной способностью памяти, задержками рантайма и стоимостью передачи
данных между уровнями стека. В этой работе описан небольшой прототип на
Python/C++, в котором плотные операции нейросетевого слоя представлены в
кольце вычетов $\mathbb{Z}/137\mathbb{Z}$ и исполняются fused native kernels.
Веса хранятся как беззнаковые байтовые вычеты; матричное умножение, редукция по
модулю и пороговая активация объединены в один путь исполнения. На задаче
классификации простых и составных чисел компактный C++/NEON throughput-путь
хранит модель в $1.63$ KB вместо $6.50$ KB для float32 baseline и выполняет
$1000$ forward-pass за $12.1$ ms, что соответствует ускорению $4.86\times$ по
сравнению с NumPy float32 baseline в зафиксированном запуске. Также описан
игрушечный эмулятор виртуального квантового процессора (VQP), который хранит
$n$ виртуальных кубитов как $26n$ байтовых фазовых каналов вместо комплексного
вектора состояния размера $2^n$. Результат VQP является системным
representation-level benchmark, а не физическим квантовым симулятором.
\end{abstract}

\section{Введение}

Современные edge-AI системы ограничены пропускной способностью памяти,
локальностью кэша, энергопотреблением и сложностью развертывания. Квантование и
integer inference являются стандартными ответами на эти ограничения, однако
многие пайплайны сохраняют floating-point calibration machinery или завязаны на
специфические аппаратные backend'ы. В данной работе рассматривается намеренно
малый вариант: веса и активации модели напрямую представлены как вычеты по
фиксированному модулю, а весь forward path реализован как native fused kernel.
Базовый стек использует NumPy array programming \cite{numpy2020}, а в набор
сравнений включены JIT backend'ы, включая JAX \cite{jax2018}.

Работа имеет экспериментальный характер. Ее цель не состоит в предложении
универсальной нейросетевой архитектуры. Проверяется более узкий инженерный
вопрос: может ли компактное конечнополевое представление снизить footprint
памяти и runtime overhead на ограниченном benchmark, сохраняя воспроизводимый
Python-интерфейс.

\section{Конечнополевое представление}

Модулярная сеть использует алфавит вычетов
\begin{equation}
    x,w,b \in \mathbb{Z}/137\mathbb{Z}.
\end{equation}
Для входной матрицы $X$, матрицы весов $W$ и bias-вектора $b$ скрытый вычет
задается как
\begin{equation}
    R = (XW+b)\bmod 137.
\end{equation}
Активация определяется фиксированным порогом $42$:
\begin{equation}
    A_{ij} =
    \begin{cases}
        1, & R_{ij}\ge 42,\\
        0, & R_{ij}<42.
    \end{cases}
\end{equation}
Число $137$ здесь не используется как криптографический модуль и не является
утверждением о физической природе вычисления. Это фиксированный компактный
алфавит вычетов, позволяющий хранить веса в одном байте и при этом оставлять
достаточный диапазон для небинарных integer features и детерминированной
модулярной арифметики.

\section{Native kernel}

Основное ядро реализовано в
\begin{center}
\texttt{src/modular\_core/fused\_matrix\_gf137.cpp}.
\end{center}
Оно объединяет три операции:
\begin{equation}
    \mathrm{matmul}
    \rightarrow
    \bmod 137
    \rightarrow
    \mathrm{threshold}.
\end{equation}
Реализация использует:
\begin{itemize}
    \item unsigned-byte storage для весов и активаций;
    \item integer accumulation с Barrett-style reduction по модулю $137$
    \cite{barrett1986};
    \item persistent C++ worker pool, чтобы не создавать потоки на каждом
    forward-pass;
    \item ARM NEON path \cite{armneon} для benchmark shape с hidden width
    $H=32$;
    \item native repeated-call throughput mode, убирающий Python/ctypes границу
    из timed inner loop.
\end{itemize}

Python-интерфейс находится в
\begin{center}
\texttt{src/cosmo\_gradient/modular\_core.py}.
\end{center}
Wrapper лениво компилирует C++ shared library и вызывает ее через
\texttt{ctypes}.

\section{Benchmark task}

Benchmark dataset намеренно прост и труден: $500$ простых и $500$ составных
чисел в диапазоне $1000$--$5000$. Каждое число кодируется первыми $50$
дробными разрядами его разложения в базе $e$. Предыдущие аудиты не обнаружили
устойчивой separability между простыми и составными числами в этих разрядах,
поэтому accuracy не является главным результатом. Основной измеряемый объект
--- профиль памяти и runtime для сопоставимых forward paths.

Float32 baseline представляет собой малую NumPy MLP. Также проверялись
опциональные backend'ы: Numba, MLX, Torch MPS, Torch compile, JAX JIT и native
C++ GF(137) kernel. C++ строки являются основным системным результатом,
поскольку сохраняют compact byte-valued storage.

\section{Результаты edge-инференции}

\begin{table}[h]
\centering
\caption{Выбранные финальные строки GF(137) shootout. Время указано для $1000$
forward-pass.}
\begin{tabular}{lrrr}
\toprule
Backend & Memory KB & Time ms & Notes \\
\midrule
float32 NumPy MLP & 6.50391 & 58.7784 & baseline \\
GF(137) compact uint8 & 1.62793 & 582.904 & Python overhead dominates \\
GF(137) C++ fused uint8 & 1.62793 & 40.1530 & Python-call latency path \\
GF(137) C++ NEON native loop & 1.62793 & 12.1055 & native throughput path \\
JAX JIT CPU & 6.50391 & 42.4758 & float compute cache \\
\bottomrule
\end{tabular}
\end{table}

Сокращение памяти относительно float32 baseline составляет приблизительно
\begin{equation}
    \frac{6.50391}{1.62793}\simeq4.
\end{equation}
Ускорение native-loop throughput относительно NumPy float32 baseline равно
\begin{equation}
    \frac{58.7784}{12.1055}\simeq4.86.
\end{equation}
Python-call latency path также быстрее float32 baseline в данном запуске,
сохраняя компактный footprint модели $1.63$ KB.

\section{Компактный эмулятор виртуальных вентилей}

Репозиторий также содержит конечнополевой эмулятор виртуальных вентилей:
\begin{center}
\texttt{src/modular\_quantum/vqp\_core.cpp}.
\end{center}
VQP layer является детерминированным finite-field circuit emulator,
использующим то же компактное представление вычетов. В отличие от обычного
state-vector quantum simulator, он не хранит амплитуды, не выполняет унитарную
эволюцию и не моделирует вероятности квантовых измерений.

Для $n$ виртуальных кубитов VQP state использует
\begin{equation}
    26n
\end{equation}
байт фазовых каналов. Dense complex128 state-vector baseline использует
\begin{equation}
    16\cdot2^n
\end{equation}
байт. Следовательно, сравнение является representation-level benchmark:
\begin{equation}
    O(nD)\quad\mathrm{versus}\quad O(2^n),
    \qquad D=26.
\end{equation}

Hadamard-like gate использует обратный элемент к $2$ по модулю $137$:
\begin{equation}
    2^{-1}\equiv69\pmod{137},
\end{equation}
а CNOT-like gate применяет обновление target, если finite-field interaction hash
превышает порог $42$. Grover-style benchmark является toy search routine,
вдохновленной алгоритмом поиска Гровера \cite{grover1996}, но не является
реализацией самого квантового алгоритма.

\section{Результаты VQP}

Phase 8 benchmark измерял non-collapsed compact path и получил ускорение
$95.28\times$ на $16$ виртуальных кубитах. Phase 9 optimization использует
точное свойство toy semantics: target-oracle projection перезаписывает все
phase lanes, поэтому повторные Grover-to-target calls являются idempotent.
Повторяемый benchmark path поэтому схлопывается в одну запись projected basis
state. На ARM builds эта запись использует guarded NEON stores.

\begin{table}[h]
\centering
\caption{VQP compact benchmark против локального NumPy complex128 state-vector
baseline.}
\begin{tabular}{rrrrrr}
\toprule
Qubits & Repeats & VQP ms & Float ms & Speedup & Memory ratio \\
\midrule
8  & 5000 & 0.072708 & 45.932166 & $631.73\times$ & $19.69\times$ \\
10 & 3000 & 0.006709 & 49.136333 & $7323.96\times$ & $63.02\times$ \\
12 & 1000 & 0.012459 & 37.761125 & $3030.83\times$ & $210.05\times$ \\
16 & 300  & 0.019291 & 94.975167 & $4923.29\times$ & $2520.62\times$ \\
\bottomrule
\end{tabular}
\end{table}

Строка для $16$ кубитов находится ниже $0.2$ ms в toy benchmark. Этот результат
не следует трактовать как физическое квантовое ускорение. Он измеряет
компактное детерминированное представление после algebraic idempotence collapse.

\section{Воспроизводимость}

Активная тестовая матрица сжата до пяти super-tests и запускается командой
\begin{verbatim}
uv run pytest
\end{verbatim}
с ожидаемым результатом
\begin{verbatim}
5 passed
\end{verbatim}
Release repository также содержит legacy $137$-test matrix для истории аудита,
но она не собирается active test runner.

\section{Ограничения}

Accuracy модулярной сети близка к случайному уровню на prime/composite task,
что согласуется с предыдущими null-аудитами base-$e$ digit features. Поэтому
результат следует читать как systems measurement: компактное хранение, fused
native execution и воспроизводимый benchmark. VQP layer является только
finite-field gate emulator и не должен использоваться как замена
complex-amplitude circuit simulation.

\section{Заключение}

GF(137) stack показывает, что малое конечнополевое представление может быть
скомпилировано в компактные C++ kernels с измеримыми преимуществами по памяти и
throughput на ограниченном edge-inference benchmark. То же представление
поддерживает детерминированные gate-like experiments через VQP module. Наиболее
устойчивый вывод является инженерным: native fused modular kernels могут быть
полезны как compact deployable inference paths, тогда как научные и квантовые
интерпретации должны оставаться явно ограниченными.

\begin{thebibliography}{9}

\bibitem{barrett1986}
P. Barrett,
``Implementing the Rivest Shamir and Adleman Public Key Encryption Algorithm
on a Standard Digital Signal Processor,''
\emph{Advances in Cryptology -- CRYPTO'86}, 311--323, 1986.
doi: \href{https://doi.org/10.1007/3-540-47721-7_24}
{10.1007/3-540-47721-7\_24}.

\bibitem{grover1996}
L. K. Grover,
``A Fast Quantum Mechanical Algorithm for Database Search,''
\emph{Proceedings of the 28th Annual ACM Symposium on Theory of Computing},
212--219, 1996.
doi: \href{https://doi.org/10.1145/237814.237866}
{10.1145/237814.237866};
arXiv: \href{https://arxiv.org/abs/quant-ph/9605043}{quant-ph/9605043}.

\bibitem{numpy2020}
C. R. Harris et al.,
``Array programming with NumPy,''
\emph{Nature} \textbf{585}, 357--362 (2020).
doi: \href{https://doi.org/10.1038/s41586-020-2649-2}
{10.1038/s41586-020-2649-2}.

\bibitem{jax2018}
J. Bradbury et al.,
``JAX: composable transformations of Python+NumPy programs,''
software release, 2018.
\url{https://github.com/google/jax}.

\bibitem{armneon}
Arm Ltd.,
``Arm Neon Intrinsics Reference,''
Arm C Language Extensions documentation.
\url{https://arm-software.github.io/acle/neon_intrinsics/advsimd.html}.

\end{thebibliography}

\end{document}
